\chapter{Stationarity and harmonic-map energy}\label{sec:harmonic-map-energy}
The domain-variation calculation is the stress-energy form of Noether's theorem in \cite{Hel01}, Chapter 3, pp.~29--33. We derive its radial consequences in any domain dimension, with the energy normalization $E(u;U)=\frac12\int_U|Du|^2$. Weak harmonicity and stationarity under domain changes are stated as separate conditions.

\begin{definition}[Weak harmonicity and domain stationarity]\label{def:stationary-weak-harmonic-map}
Let $\Omega\subset\mathbb R^m$ be open and $u\in W^{1,2}_{\mathrm{loc}}(\Omega;\mathbb R^N)$ with $|u|=1$ almost everywhere. It is weakly harmonic if
\[
 \int_\Omega Du:D\zeta=\int_\Omega |Du|^2u\cdot\zeta
 \quad(\zeta\in C_c^\infty(\Omega;\mathbb R^N)).
\]
It is stationary under domain variations if, in addition,
\[
 \int_\Omega\left(\frac12|Du|^2\operatorname{div}X
       -\sum_{i,j}\langle u_i,u_j\rangle\partial_iX_j\right)=0
 \quad(X\in C_c^\infty(\Omega;\mathbb R^m)).
\]
We call a map satisfying both conditions a stationary weak harmonic map. The second identity is an additional hypothesis in the weak theory.
\end{definition}

\begin{proposition}[The stress tensor and inner variations]\label{prop:harmonic-map-stress-tensor}
For a smooth harmonic map $u:\Omega\to S^{N-1}$, set
\[
 S_{ij}=\tfrac12|Du|^2\delta_{ij}-\langle u_i,u_j\rangle.
\]
Then $\sum_i\partial_iS_{ij}=0$. Moreover, for the flow $\Phi_t$ of a compactly supported domain field $X$, the derivative of the energy of $u\circ\Phi_t^{-1}$ is $\int S_{ij}\partial_iX_j$. In particular every smooth harmonic map is stationary under domain variations. Compare \cite{Hel01}, Theorem 3.2; its Hamiltonian tensor is $-S$.
\end{proposition}
\begin{proof}
Differentiating and canceling mixed-derivative terms gives
\[
 \sum_i\partial_iS_{ij}=-\langle\Delta u,u_j\rangle=|Du|^2\langle u,u_j\rangle=0.
\]
For the variation, change variables from $\Phi_t(x)$ to $x$ on a relatively compact domain containing the support. The Jacobian has derivative $\operatorname{div}X$, while $D\Phi_t^{-1}$ has derivative $-DX$. Differentiating $\frac12|Du(D\Phi_t)^{-1}|^2\det D\Phi_t$ gives exactly the claimed integrand. Integrating the divergence-free tensor against $DX$ gives zero by integration by parts. This smooth computation does not justify using $Du\,X$ as an admissible test in a merely weak equation; that is why Definition \ref{def:stationary-weak-harmonic-map} retains domain stationarity separately.
\end{proof}

\begin{theorem}[Radial energy monotonicity]\label{thm:stationary-harmonic-energy-monotonicity}
Let $m\ge2$, let $u$ be a stationary weak harmonic map, and assume $\overline B_R(a)\subset\Omega$. Define
\[
 \Theta_u(a,r)=r^{2-m}\frac12\int_{B_r(a)}|Du|^2.
\]
For $0<\rho<R$,
\[
 \Theta_u(a,R)-\Theta_u(a,\rho)
 =\int_{B_R(a)\setminus B_\rho(a)}|x-a|^{2-m}|\partial_r u|^2\,dx\ge0,
\]
where $\partial_r u=Du(x)(x-a)/|x-a|$. In particular $\Theta_u(a,r)$ is nondecreasing and has a finite nonnegative limit as $r\downarrow0$. The proof derives the radial identity from the domain-variation framework of \cite{Hel01}, Theorem 3.2, rather than importing weak regularity results.
\end{theorem}
\begin{proof}
Translate $a$ to zero and write $e=|Du|^2/2$ and $E(r)=\int_{B_r}e$. In the domain-stationarity identity use $X(x)=\eta(|x|)x$, with $\eta$ smooth and constant near zero. Since $\operatorname{tr}S=(m-2)e$ and $S(x,x)/|x|=|x|(e-|\partial_ru|^2)$, this yields
\[
 0=\int\big((m-2)\eta e+r\eta'(e-|\partial_ru|^2)\big).
\]
Approximate the indicator of $[0,s]$ by such cutoffs. Polar integration and one-dimensional Lebesgue differentiation give, for almost every $s$,
\[
 (m-2)E(s)=s\int_{\partial B_s}(e-|\partial_ru|^2).
\]
Here all surface integrals are the almost-everywhere slices given by polar integration of integrable functions. Since $E'(s)=\int_{\partial B_s}e$ almost everywhere,
\[
 \frac{d}{ds}(s^{2-m}E(s))=s^{2-m}\int_{\partial B_s}|\partial_ru|^2.
\]
Both sides are integrable on every closed radius interval away from zero, and $E$ is absolutely continuous there. Integrating and applying polar integration proves the formula for all $0<\rho<R$. Monotonicity and the finite value at $R$ give the finite limit at zero. The argument uses domain stationarity; it does not need to differentiate the weak harmonic-map equation.
\end{proof}

\begin{proposition}[Equality, homogeneity and smooth-point density]\label{prop:harmonic-energy-equality}
Equality of $\Theta_u(a,R)$ and $\Theta_u(a,\rho)$ in Theorem \ref{thm:stationary-harmonic-energy-monotonicity} is equivalent to $\partial_ru=0$ almost everywhere on the annulus. Equivalently $u(a+r\omega)$ is independent of $r\in(\rho,R)$ for almost every direction $\omega$, up to Sobolev representatives. At a point where $u$ is smooth, $\lim_{r\downarrow0}\Theta_u(a,r)=0$.
\end{proposition}
\begin{proof}
The integral in the monotonicity identity is of a nonnegative function with a strictly positive weight. It is zero precisely when the radial derivative vanishes almost everywhere. On any annulus bounded away from its center, polar charts are smooth local changes of variables. The absolute continuity on almost every coordinate line of a Sobolev function, applied in these charts and then exhausted by compact subannuli, makes vanishing radial derivative equivalent to radial constancy. At a smooth point the energy density is bounded on a small ball, so $\Theta_u(a,r)\le Cr^{2-m}|B_r|=C\omega_m r^2\to0$.
\end{proof}

\begin{corollary}[Finite total energy in dimensions above two]\label{cor:finite-energy-harmonic-liouville}
If $m\ge3$ and $u:\mathbb R^m\to S^{N-1}$ is a stationary weak harmonic map with $\int_{\mathbb R^m}|Du|^2<\infty$, then $u$ is constant almost everywhere. The same conclusion holds for a smooth harmonic map under this energy hypothesis. Dimension two admits nonconstant finite-energy examples.
\end{corollary}
\begin{proof}
For any fixed $r>0$, monotonicity gives $0\le\Theta_u(0,r)\le\Theta_u(0,R)\le \frac12R^{2-m}\int_{\mathbb R^m}|Du|^2\to0$. Thus $Du=0$ almost everywhere on each centered ball and hence on all of $\mathbb R^m$. The standard line characterization of Sobolev functions on connected domains gives an almost-everywhere constant representative. Smooth harmonic maps satisfy stationarity by Proposition \ref{prop:harmonic-map-stress-tensor}. Example \ref{ex:stereographic-energy-concentration} below gives the dimension-two contrast.
\end{proof}
